% Source: https://artofproblemsolving.com/community/c3h1687394_the_n_ns
% Stable identity: aops-topic-1687394-q01
Write $\sqrt{\ }^{\,k}$ for $k$ repeated square roots, so $\sqrt{\ }^{\,k}\,n=n^{\dfrac{1}{2^k}}$. Since $n\ge4>1$, logarithms to base $n$, $\sqrt n$ and $\dfrac{1}{2}$ make sense. We use the building blocks
\[
\log_{\sqrt n}n=2,\qquad \log_n\sqrt n=\dfrac{1}{2},\qquad \log_n\left(\sqrt{\ }^{\,k}\,n\right)=\dfrac{1}{2^k},\qquad \{n\}=0,
\]
which use two, two, two and one copies of $n$ respectively.

\textbf{Negative integers.} For $k\ge1$, using four copies of $n$,
\[
\log_{\log_{\sqrt n}n}\Bigl(\log_n\bigl(\sqrt{\ }^{\,k}\,n\bigr)\Bigr)=\log_2\dfrac{1}{2^k}=-k .
\]

\textbf{Positive integers.} For $k\ge1$, again with four copies,
\[
\log_{\log_n\sqrt n}\Bigl(\log_n\bigl(\sqrt{\ }^{\,k}\,n\bigr)\Bigr)=\log_{\dfrac{1}{2}}\dfrac{1}{2^k}=k .
\]

\textbf{Zero.} $(n-n)\cdot(n-n)=0$ uses four copies.

\textbf{Using all $n$ copies.} Each of these expressions uses exactly four copies of $n$. When $n>4$, add $\{n\}=0$ once for each of the remaining $n-4$ copies; this does not change the value. Hence every integer can be produced with exactly $n$ copies of $n$. $\blacksquare$

(This is the idea behind Dirac's solution of the classical ``four $2$s'' puzzle.)
